\documentclass[11pt]{article}
\usepackage[margin=1in]{geometry}
\usepackage{amsmath,amssymb}
\usepackage{bm}
\usepackage{booktabs}
\usepackage{xcolor}
\usepackage{hyperref}
\hypersetup{colorlinks=true,linkcolor=blue,urlcolor=blue}

\newcommand{\rhat}{\hat{\mathbf r}}
\newcommand{\zhat}{\hat{\mathbf z}}
\newcommand{\that}{\hat{\boldsymbol\theta}}
\newcommand{\dair}{d_{\mathrm{air}}}
\newcommand{\Emodem}{E_{\mathrm{ModEM}}}
\newcommand{\Eg}{E_{\mathrm{global1d}}}

\title{The toroidal source potential, the uniform inducing field,\\
and the \textsc{global1d}$\leftrightarrow$\textsc{ModEM} normalization factor}
\author{Anna Kelbert with Claude}
\date{July 2026}

\begin{document}
\maketitle

\begin{abstract}
\textsc{global1d} (\texttt{field1d.f90}/\texttt{field1d\_s2.f90}) drives its solution with an
\emph{external toroidal potential amplitude} $\mathrm{coeff}(l,m)$; \textsc{ModEM}'s regional 3-D
solver drives its solution with a \emph{boundary electric field} of unit amplitude. Because a
\emph{potential} and a \emph{field} are related by Faraday's law $\nabla\times\mathbf E=i\omega\mu_0
\mathbf H$, the raw electric fields the two codes produce for their respective unit sources differ by
a single multiplicative complex constant per period, $c=\Eg/\Emodem=i\,\omega\,\mu_0\,G$, with $G$ a
real geometric constant set by the \emph{air-layer thickness} and the surface magnetic response. This
note explains the physics (the toroidal potential defines a quasi-uniform \emph{magnetic} field, not
an electric one), derives $G$ from the air-layer geometry, and defines a normalized output preset
that matches \textsc{ModEM}'s normalization. There is \textbf{no bug and no sign/time-convention
error}: $c$ cancels identically in the impedance $Z=E/H$, so the two codes agree on $Z$, apparent
resistivity and transfer functions to $<1\%$. The apparent ``$-\mathrm{Re},+\mathrm{Im}$ vs
$+\mathrm{Re},+\mathrm{Im}$ sign flip'' in raw fields is simply the ${\approx}90^\circ$ phase of the
$i\omega\mu_0$ factor.
\end{abstract}

\section{The toroidal potential and the fields it produces}
Under the quasi-static approximation a divergence-free field on a sphere admits a Mie
(poloidal--toroidal) decomposition. In the insulating air there are no currents, so the poloidal
electric potential $P$ vanishes (Sun \& Egbert 2012, eqs.~21--22) and an external MT/ionospheric
source drives only the toroidal potential $T$. With $P\equiv0$, eqs.~(5)--(6) reduce to
\begin{equation}
\mathbf E=-i\omega\mu_0\,\frac{\rhat}{r}\times\nabla_a T,\qquad
\mathbf H=\frac1r\nabla_a\frac{\partial T}{\partial r}-\frac{\rhat}{r^2}\nabla_a^2 T .
\end{equation}
Writing $T(r,\theta,\phi)=\sum_{l,m}\mathrm{coeff}(l,m)\,T_l(r)\,Y_l^m(\theta,\phi)$ and using
$\rhat\times\that=\hat{\boldsymbol\phi}$, $\nabla_a^2 Y_l^m=-l(l+1)Y_l^m$, the components are
(\texttt{field1d\_s2.f90}, l.\,42--47)
\begin{equation}
\begin{aligned}
E_\theta&=+\tfrac{i\omega\mu_0}{r}T_l Y_p, & E_\phi&=-\tfrac{i\omega\mu_0}{r}T_l Y_t, & E_r&=0,\\
H_r&=+\tfrac{l(l+1)}{r^2}T_l Y, & H_\theta&=+\tfrac1r T_l' Y_t, & H_\phi&=+\tfrac1r T_l' Y_p,
\end{aligned}
\end{equation}
with $Y_t=\partial Y/\partial\theta$, $Y_p=(1/\sin\theta)\partial Y/\partial\phi$. The single scalar
$T$ generates a \textbf{toroidal electric field} ($E_r=0$; the name refers to this) and a
\textbf{poloidal magnetic field} ($H_r\neq0$, the vertical/$Z$ component at the ground)---the
standard deep-Earth induction picture.

\section{The (external) toroidal potential is a quasi-uniform magnetic field}
In the air $T$ is harmonic radially, $T_l(r)=\alpha r^{l+1}+\beta r^{-l}$: an external ($\alpha$,
source) plus internal ($\beta$, induced) part. \textsc{global1d} sets $\mathrm{coeff}(l,m)=1$
$\Rightarrow\alpha=1$ (unit coefficient of $r^{l+1}$; \texttt{field1d\_s2.f90}, l.\,437--458). For
$l=1,m=0$ with $Y_1^0=\sqrt{3/4\pi}\,\cos\theta$ and $T_1^{\mathrm{ext}}=\alpha r^2$,
\begin{equation}
H_r=2\alpha\sqrt{\tfrac{3}{4\pi}}\cos\theta,\qquad
H_\theta=-2\alpha\sqrt{\tfrac{3}{4\pi}}\sin\theta,
\end{equation}
and since $\cos\theta\,\rhat-\sin\theta\,\that=\zhat$,
\begin{equation}
\boxed{\ \mathbf H^{\mathrm{ext}}=2\alpha\sqrt{\tfrac{3}{4\pi}}\,\zhat\ }
\qquad\text{a spatially uniform axial magnetic field.}
\end{equation}
The induced electric field is \emph{not} uniform and carries the explicit $i\omega\mu_0$:
$E_\phi=+i\omega\mu_0\,\alpha\,r\sqrt{3/4\pi}\,\sin\theta$. \textbf{Thus the (external, $l=1$)
toroidal potential defines a quasi-uniform inducing \emph{magnetic} field, and $E$ is the
$i\omega\mu_0$-scaled induced response.}

\section{\textsc{global1d} vs \textsc{ModEM}: $c=i\omega\mu_0 G$}
\textsc{global1d}'s unit source ($\propto$ uniform $H$) and \textsc{ModEM}'s unit source (a boundary
$E$-field) differ by the potential-vs-field relationship. The measured ratio $c=\Eg/\Emodem$ is
(i)~\emph{constant over depth}---a multiplicative, not antilinear, relationship;
(ii)~\emph{scaling as $\omega$}: $|c|/\omega$ constant and $\arg c\to90^\circ$ at short period
$\Rightarrow c=i\omega\mu_0 G$; (iii)~\emph{mode-independent}. The $i$ is Faraday's law; the small
residual ($\arg c$ drifting $90^\circ\!\to\!95^\circ$ over $T=10\to1000\,$s) is the earth's induction
response (\S4) plus \textsc{ModEM}'s flat-vs-spherical boundary approximation.

\section{Derivation of $G$ from the air-layer geometry}
\paragraph{ModEM's surface field.} \textsc{ModEM} drives a unit $E$ at the top of the domain, a
height $\dair\approx1000\,$km above the surface, and solves a 1-D column
(\texttt{Mod2d/WSfwd1Dmod.f90}). In the insulating air $d^2E/dz^2=0$, so $E$ is linear in height and
$H$ is uniform (as in \S2). With $E(0)=1$ and surface impedance $Z_e=E_s/H_s$,
\begin{equation}
\Emodem(\text{surface})=\frac{Z_e}{Z_e-i\omega\mu_0\dair}\approx\frac{-Z_e}{i\omega\mu_0\dair},
\end{equation}
since $|Z_e|\ll\omega\mu_0\dair$ for a $1000\,$km air column. \textsc{ModEM}'s surface field is set
by the air thickness.

\paragraph{The ratio.} \textsc{global1d} obeys $E_\phi=Z_e H_\theta$ at the surface (same $Z_e$).
Taking the ratio, $Z_e$ cancels:
\begin{equation}
c=\frac{\Eg}{\Emodem}=H_\theta(Z_e-i\omega\mu_0\dair)=E_\phi-i\omega\mu_0\dair H_\theta
=-i\omega\mu_0\,H_\theta\,(\dair+C),
\end{equation}
where $C\equiv Z_e/(i\omega\mu_0)$ is the C-response ($\approx4,28,132\,$km at $T=10,100,1000\,$s).
Hence
\begin{equation}
\boxed{\ G=-H_\theta(\text{surface})\,(\dair+C)\;\approx\;-H_\theta\,\dair\ }
\end{equation}
the surface horizontal field times the air-layer thickness; $C$ is the period-dependent residual.

\paragraph{The geometric value of $H_\theta$.} From $T_1(r)=r^2+\beta/r$ at $r_0$,
$T_1'(r_0)=2r_0-\beta/r_0^2$, so
\begin{equation}
H_\theta(\text{surface})=-\Big(2-\tfrac{\beta}{r_0^3}\Big)\sqrt{\tfrac{3}{4\pi}}\sin\theta
=-(2-Q_1)\sqrt{\tfrac{3}{4\pi}}\sin\theta,
\end{equation}
with $Q_1=\beta/r_0^3$ the $l=1$ Q-response (limits $Q_1=0$ insulator, $Q_1=l/(l+1)=\tfrac12$ perfect
conductor). At MT/geomagnetic periods the mantle is a good conductor, $Q_1\to\tfrac12$, and at the
equator
\begin{equation}
\boxed{\ G=\tfrac32\sqrt{\dfrac{3}{4\pi}}\ \dair\;\approx\;0.733\,\dair\ }
\end{equation}
Here $\tfrac32=(2-Q_1)$ combines external ($2$) and induced ($-Q_1$) parts, $\sqrt{3/4\pi}$ is the
$Y_1^0$ amplitude, and $\dair=\texttt{grid\%r(1)}\!\cdot\!1000-r_0\approx1.0\times10^{6}\,\mathrm{m}$.

\paragraph{Numerical verification.}
\begin{center}
\begin{tabular}{cccc}
\toprule
$T$ (s) & $|H_\theta|$ & $|G|/\dair$ (meas.) & $\arg G$\\
\midrule
10 & 0.718 & 0.719 & $+0.65^\circ$\\
100 & 0.721 & 0.730 & $+1.59^\circ$\\
1000 & 0.720 & 0.783 & $+4.73^\circ$\\
\bottomrule
\end{tabular}
\end{center}
The derived $\tfrac32\sqrt{3/4\pi}=0.733$ straddles the measured $|G|/\dair$ (best at $100\,$s); the
drift and growing $\arg G$ are the $(\dair+C)$ correction. $H_\theta$ is nearly period-independent,
confirming it is geometric in the deep-induction limit ($Q_1\approx\tfrac12$).

\section{The normalized output preset}
To express \textsc{global1d}'s fields in \textsc{ModEM}'s normalization, multiply by $1/c$:
\begin{equation}
(\mathbf E,\mathbf H)_{\mathrm{norm}}=\frac{1}{i\omega\mu_0 G}(\mathbf E,\mathbf H),
\qquad G=\tfrac32\sqrt{\tfrac{3}{4\pi}}\,\dair,
\end{equation}
with $\dair$ from the grid and $\omega$ from the period. Applying the same complex factor to $E$ and
$H$ (equivalently rotating $\mathrm{coeff}$ by $1/i=-i$ and scaling) \textbf{preserves $Z=E/H$}: the
physics is untouched; only the source-normalization convention changes. Implemented as the
\texttt{EGBERTKELBERT2012\_MODEM} preset (\texttt{output\_convention.f90}). The factor removes the
exact $i\omega\mu_0$ and the air-thickness scaling; residuals (the $C$-response, finite $Q$, and
flat-vs-spherical) leave a few-percent amplitude and ${\lesssim}5^\circ$ phase mismatch at long
period, $\to 0$ at short period---genuine physics differences, not artifacts.

\section{Validation against Cartesian ModEM}
The spherical \textsc{ModEM} runs used to \emph{measure} $c$ (\texttt{Mod3DMT\_SP2\_SPH}) require
building with \texttt{-DBUILD\_SPHERICAL -DFORCE\_SPHERICAL}, bypassing a guard that warns the
spherical version is \emph{experimental and ``will not work with traditional MT source setup''}---and
the internal \texttt{COMPUTE\_BC} boundary condition IS that traditional (WS/REBOCC per-column) setup.
The derivation was therefore cross-checked against \textbf{Cartesian ModEM} (trusted and accurate) on
a $\rho{=}100\,\Omega$m halfspace with the same $1000\,$km/12-layer air.

\emph{Air-column mechanism} (the heart of $G$): Cartesian \textsc{ModEM}'s surface $E$ equals the
analytic $Z_e/(Z_e-i\omega\mu_0\dair)$ (with $Z_e=-i\omega\mu_0/\sqrt{-i\omega\mu_0\sigma}$) to
$0.2$--$0.6\%$, $\pm0.1^\circ$ at every period from $4$ to $4000\,$s.

\emph{Full preset}, $r=\Eg^{\mathrm{norm}}(\text{equator})/\Emodem^{\mathrm{cart}}$:
\begin{center}
\begin{tabular}{cccc}
\toprule
$T$ (s) & 3.98 & 1000 & 3981\\
\midrule
$|r|$ & 1.006 & 1.073 & 1.145\\
$\arg r$ & $+0.2^\circ$ & $+3.6^\circ$ & $+6.4^\circ$\\
\bottomrule
\end{tabular}
\end{center}
The match is essentially perfect in the flat-earth limit ($|r|=1.006$, $+0.2^\circ$ at $4\,$s, where
$G$ is exact and $Q_1\to\tfrac12$) and the residual grows monotonically and cleanly with period as
curvature enters; Cartesian \textsc{ModEM} is laterally uniform to $0.00\%$. This validates the
$i\omega\mu_0G$ framework on trusted code and shows the residual is the \textbf{genuine
flat-vs-spherical} physics---\textsc{global1d} being the more accurate spherical solution---not an
artifact of the experimental spherical BC.

\paragraph{Generality.} The derivation is for $l=1$ (P10 / MT modes) at the equator in the
deep-induction limit. In general $H_\theta=-[(l+1)-l\,Q_l]\,r_0^{\,l-1}\,\partial_\theta Y_l^0$
depends on $(l,\theta)$, so a single scalar $G$ is exact only for the $l=1$ equatorial case the MT
modes reduce to---the relevant limit for MT interoperability.

\end{document}
